% Options for packages loaded elsewhere
\PassOptionsToPackage{unicode}{hyperref}
\PassOptionsToPackage{hyphens}{url}
\PassOptionsToPackage{dvipsnames,svgnames,x11names}{xcolor}
%
\documentclass[
  a4paper,
  DIV=11,
  numbers=noendperiod]{scrartcl}

\usepackage{amsmath,amssymb}
\usepackage{iftex}
\ifPDFTeX
  \usepackage[T1]{fontenc}
  \usepackage[utf8]{inputenc}
  \usepackage{textcomp} % provide euro and other symbols
\else % if luatex or xetex
  \usepackage{unicode-math}
  \defaultfontfeatures{Scale=MatchLowercase}
  \defaultfontfeatures[\rmfamily]{Ligatures=TeX,Scale=1}
\fi
\usepackage{lmodern}
\ifPDFTeX\else  
    % xetex/luatex font selection
\fi
% Use upquote if available, for straight quotes in verbatim environments
\IfFileExists{upquote.sty}{\usepackage{upquote}}{}
\IfFileExists{microtype.sty}{% use microtype if available
  \usepackage[]{microtype}
  \UseMicrotypeSet[protrusion]{basicmath} % disable protrusion for tt fonts
}{}
\makeatletter
\@ifundefined{KOMAClassName}{% if non-KOMA class
  \IfFileExists{parskip.sty}{%
    \usepackage{parskip}
  }{% else
    \setlength{\parindent}{0pt}
    \setlength{\parskip}{6pt plus 2pt minus 1pt}}
}{% if KOMA class
  \KOMAoptions{parskip=half}}
\makeatother
\usepackage{xcolor}
\usepackage[landscape]{geometry}
\setlength{\emergencystretch}{3em} % prevent overfull lines
\setcounter{secnumdepth}{-\maxdimen} % remove section numbering
% Make \paragraph and \subparagraph free-standing
\makeatletter
\ifx\paragraph\undefined\else
  \let\oldparagraph\paragraph
  \renewcommand{\paragraph}{
    \@ifstar
      \xxxParagraphStar
      \xxxParagraphNoStar
  }
  \newcommand{\xxxParagraphStar}[1]{\oldparagraph*{#1}\mbox{}}
  \newcommand{\xxxParagraphNoStar}[1]{\oldparagraph{#1}\mbox{}}
\fi
\ifx\subparagraph\undefined\else
  \let\oldsubparagraph\subparagraph
  \renewcommand{\subparagraph}{
    \@ifstar
      \xxxSubParagraphStar
      \xxxSubParagraphNoStar
  }
  \newcommand{\xxxSubParagraphStar}[1]{\oldsubparagraph*{#1}\mbox{}}
  \newcommand{\xxxSubParagraphNoStar}[1]{\oldsubparagraph{#1}\mbox{}}
\fi
\makeatother

\usepackage{color}
\usepackage{fancyvrb}
\newcommand{\VerbBar}{|}
\newcommand{\VERB}{\Verb[commandchars=\\\{\}]}
\DefineVerbatimEnvironment{Highlighting}{Verbatim}{commandchars=\\\{\}}
% Add ',fontsize=\small' for more characters per line
\usepackage{framed}
\definecolor{shadecolor}{RGB}{241,243,245}
\newenvironment{Shaded}{\begin{snugshade}}{\end{snugshade}}
\newcommand{\AlertTok}[1]{\textcolor[rgb]{0.68,0.00,0.00}{#1}}
\newcommand{\AnnotationTok}[1]{\textcolor[rgb]{0.37,0.37,0.37}{#1}}
\newcommand{\AttributeTok}[1]{\textcolor[rgb]{0.40,0.45,0.13}{#1}}
\newcommand{\BaseNTok}[1]{\textcolor[rgb]{0.68,0.00,0.00}{#1}}
\newcommand{\BuiltInTok}[1]{\textcolor[rgb]{0.00,0.23,0.31}{#1}}
\newcommand{\CharTok}[1]{\textcolor[rgb]{0.13,0.47,0.30}{#1}}
\newcommand{\CommentTok}[1]{\textcolor[rgb]{0.37,0.37,0.37}{#1}}
\newcommand{\CommentVarTok}[1]{\textcolor[rgb]{0.37,0.37,0.37}{\textit{#1}}}
\newcommand{\ConstantTok}[1]{\textcolor[rgb]{0.56,0.35,0.01}{#1}}
\newcommand{\ControlFlowTok}[1]{\textcolor[rgb]{0.00,0.23,0.31}{\textbf{#1}}}
\newcommand{\DataTypeTok}[1]{\textcolor[rgb]{0.68,0.00,0.00}{#1}}
\newcommand{\DecValTok}[1]{\textcolor[rgb]{0.68,0.00,0.00}{#1}}
\newcommand{\DocumentationTok}[1]{\textcolor[rgb]{0.37,0.37,0.37}{\textit{#1}}}
\newcommand{\ErrorTok}[1]{\textcolor[rgb]{0.68,0.00,0.00}{#1}}
\newcommand{\ExtensionTok}[1]{\textcolor[rgb]{0.00,0.23,0.31}{#1}}
\newcommand{\FloatTok}[1]{\textcolor[rgb]{0.68,0.00,0.00}{#1}}
\newcommand{\FunctionTok}[1]{\textcolor[rgb]{0.28,0.35,0.67}{#1}}
\newcommand{\ImportTok}[1]{\textcolor[rgb]{0.00,0.46,0.62}{#1}}
\newcommand{\InformationTok}[1]{\textcolor[rgb]{0.37,0.37,0.37}{#1}}
\newcommand{\KeywordTok}[1]{\textcolor[rgb]{0.00,0.23,0.31}{\textbf{#1}}}
\newcommand{\NormalTok}[1]{\textcolor[rgb]{0.00,0.23,0.31}{#1}}
\newcommand{\OperatorTok}[1]{\textcolor[rgb]{0.37,0.37,0.37}{#1}}
\newcommand{\OtherTok}[1]{\textcolor[rgb]{0.00,0.23,0.31}{#1}}
\newcommand{\PreprocessorTok}[1]{\textcolor[rgb]{0.68,0.00,0.00}{#1}}
\newcommand{\RegionMarkerTok}[1]{\textcolor[rgb]{0.00,0.23,0.31}{#1}}
\newcommand{\SpecialCharTok}[1]{\textcolor[rgb]{0.37,0.37,0.37}{#1}}
\newcommand{\SpecialStringTok}[1]{\textcolor[rgb]{0.13,0.47,0.30}{#1}}
\newcommand{\StringTok}[1]{\textcolor[rgb]{0.13,0.47,0.30}{#1}}
\newcommand{\VariableTok}[1]{\textcolor[rgb]{0.07,0.07,0.07}{#1}}
\newcommand{\VerbatimStringTok}[1]{\textcolor[rgb]{0.13,0.47,0.30}{#1}}
\newcommand{\WarningTok}[1]{\textcolor[rgb]{0.37,0.37,0.37}{\textit{#1}}}

\providecommand{\tightlist}{%
  \setlength{\itemsep}{0pt}\setlength{\parskip}{0pt}}\usepackage{longtable,booktabs,array}
\usepackage{calc} % for calculating minipage widths
% Correct order of tables after \paragraph or \subparagraph
\usepackage{etoolbox}
\makeatletter
\patchcmd\longtable{\par}{\if@noskipsec\mbox{}\fi\par}{}{}
\makeatother
% Allow footnotes in longtable head/foot
\IfFileExists{footnotehyper.sty}{\usepackage{footnotehyper}}{\usepackage{footnote}}
\makesavenoteenv{longtable}
\usepackage{graphicx}
\makeatletter
\def\maxwidth{\ifdim\Gin@nat@width>\linewidth\linewidth\else\Gin@nat@width\fi}
\def\maxheight{\ifdim\Gin@nat@height>\textheight\textheight\else\Gin@nat@height\fi}
\makeatother
% Scale images if necessary, so that they will not overflow the page
% margins by default, and it is still possible to overwrite the defaults
% using explicit options in \includegraphics[width, height, ...]{}
\setkeys{Gin}{width=\maxwidth,height=\maxheight,keepaspectratio}
% Set default figure placement to htbp
\makeatletter
\def\fps@figure{htbp}
\makeatother

<script src="https://cdnjs.cloudflare.com/ajax/libs/jspdf/2.5.1/jspdf.umd.min.js"></script>
\KOMAoption{captions}{tableheading}
\makeatletter
\@ifpackageloaded{caption}{}{\usepackage{caption}}
\AtBeginDocument{%
\ifdefined\contentsname
  \renewcommand*\contentsname{Table of contents}
\else
  \newcommand\contentsname{Table of contents}
\fi
\ifdefined\listfigurename
  \renewcommand*\listfigurename{List of Figures}
\else
  \newcommand\listfigurename{List of Figures}
\fi
\ifdefined\listtablename
  \renewcommand*\listtablename{List of Tables}
\else
  \newcommand\listtablename{List of Tables}
\fi
\ifdefined\figurename
  \renewcommand*\figurename{Figure}
\else
  \newcommand\figurename{Figure}
\fi
\ifdefined\tablename
  \renewcommand*\tablename{Table}
\else
  \newcommand\tablename{Table}
\fi
}
\@ifpackageloaded{float}{}{\usepackage{float}}
\floatstyle{ruled}
\@ifundefined{c@chapter}{\newfloat{codelisting}{h}{lop}}{\newfloat{codelisting}{h}{lop}[chapter]}
\floatname{codelisting}{Listing}
\newcommand*\listoflistings{\listof{codelisting}{List of Listings}}
\makeatother
\makeatletter
\makeatother
\makeatletter
\@ifpackageloaded{caption}{}{\usepackage{caption}}
\@ifpackageloaded{subcaption}{}{\usepackage{subcaption}}
\makeatother

\ifLuaTeX
  \usepackage{selnolig}  % disable illegal ligatures
\fi
\usepackage{bookmark}

\IfFileExists{xurl.sty}{\usepackage{xurl}}{} % add URL line breaks if available
\urlstyle{same} % disable monospaced font for URLs
\hypersetup{
  colorlinks=true,
  linkcolor={blue},
  filecolor={Maroon},
  citecolor={Blue},
  urlcolor={Blue},
  pdfcreator={LaTeX via pandoc}}


\author{}
\date{}

\begin{document}


\begin{Shaded}
\begin{Highlighting}[]
\NormalTok{//| echo: false}
\NormalTok{html\textasciigrave{}}
\NormalTok{\textless{}style\textgreater{}}
\NormalTok{.header{-}flex \{}
\NormalTok{  display: flex;}
\NormalTok{  align{-}items: center;}
\NormalTok{  margin{-}bottom: 0em;}
\NormalTok{\}}
\NormalTok{.header{-}title \{}
\NormalTok{  font{-}size: 1.5em;}
\NormalTok{  font{-}weight: bold;}
\NormalTok{\}}
\NormalTok{\textless{}/style\textgreater{}}
\NormalTok{\textasciigrave{}}
\end{Highlighting}
\end{Shaded}

Ammonia
kalkylator

\begin{Shaded}
\begin{Highlighting}[]
\NormalTok{//| echo: false}
\NormalTok{html\textasciigrave{}}
\NormalTok{\textless{}details\textgreater{}}
\NormalTok{ \textless{}summary\textgreater{}\textless{}span style="font{-}size:24px;font{-}weight:bold;color:\#03865a;"\textgreater{}Instruktion\textless{}/span\textgreater{}\textless{}/summary\textgreater{}}

\NormalTok{Detta verktyg beräknar behandlingstiden för en 5{-}log10 reduktion av Escherichia coli baserat på ureados och temperatur i behandlingsbrunn/reaktor. E. coli är en mag{-}tarm bakterie som alltid förekommer i klosettvatten och inaktiveringen av denna är ett bra mått på inaktiveringen av bakteriella smittämnen. En 5{-}log10 reduktion innebär att bakterien inaktiverats med 99,999\%. Vid en sådan inaktivering av E.coli bör klosettvattnet vara fritt från salmonella vilken förekommer i mycket lägre startkoncentrationer, om den förekommer. Saknar du startvärden kan du ange månad i rulllist och få en förväntad minimitemperatur i brunnen. Tid för nedbrytning av urea är inkluderad i behandingstiden och baseras på angiven temperatur och förutsätter omblandning.  Dessa uppskattningar är tänkta att underlätta planering av ett system för kretslopp från klosettvatten. Inför gödsling bör innehållet av E. coli och salmonella analyseras.}

\NormalTok{\textless{}/details\textgreater{}}
\NormalTok{\textless{}hr style="border{-}top:5px  \#03865a;"\textgreater{}}
\NormalTok{\textasciigrave{}}
\end{Highlighting}
\end{Shaded}

\begin{figure}

\begin{minipage}{0.50\linewidth}

\begin{Shaded}
\begin{Highlighting}[]
\NormalTok{//| echo: false}
\NormalTok{viewof dry2 = Inputs.range([0, 6], \{}
\NormalTok{  value: 0.2,}
\NormalTok{  step: 0.10,}
\NormalTok{  label: "Torrsubstans (\%)",}
\NormalTok{  style: \textasciigrave{}}
\NormalTok{    width: 50\%;}
\NormalTok{    height: 155px;}
\NormalTok{    background: linear{-}gradient(90deg, \#667eea 0\%, \#764ba2 100\%);}
\NormalTok{    border{-}radius: 8px;}
\NormalTok{    outline: none;}
\NormalTok{    {-}webkit{-}appearance: none;}
\NormalTok{    appearance: none;}
\NormalTok{  \textasciigrave{}}
\NormalTok{\})}




\NormalTok{// Input numbers}
\NormalTok{html\textasciigrave{}\textless{}div\textgreater{}}
\NormalTok{  \textless{}h4\textgreater{}\textless{}strong\textgreater{}Klosettvatten startvärden:\textless{}/strong\textgreater{}\textless{}/h4\textgreater{}}
\NormalTok{\textless{}/div\textgreater{}\textasciigrave{}}
\NormalTok{viewof dry = Inputs.range([0, 6], \{value: 0.2,step: 0.10,label: "Torrsubstans (\%)"\})}
\NormalTok{viewof kvave = Inputs.range([0, 10], \{value: 0.35,step: 0.1,label: "Kväve (g/L)"\})}
\NormalTok{viewof ammokvave = Inputs.range([0, 10], \{value: 0.3,step: 0.1,label: "Ammoniak{-}kväve (g/L)"\})}

\NormalTok{html\textasciigrave{}\textless{}div\textgreater{}}
\NormalTok{  \textless{}h4\textgreater{}Behandling:\textless{}/h4\textgreater{}}
\NormalTok{\textless{}/div\textgreater{}\textasciigrave{}}
\NormalTok{viewof va\_vol = Inputs.range([0, 5000], \{value: 1000,step: 100,label: "Klosettvattenvolym (m³)"\})}

\NormalTok{viewof urea = Inputs.range([0, 5], \{value: 1,step: 0.1,label: "Urea{-}dos (\%)"\})}

\NormalTok{options = [}
\NormalTok{\{label: "Okt{-}mars", value: 10\},}
\NormalTok{\{label: "April{-}Juni", value: 15\},}
\NormalTok{\{label: "Juli{-}Sept", value: 20\}}
\NormalTok{]}

\NormalTok{viewof selectedOption = Inputs.select(options, \{}
\NormalTok{label: "Månad (välj i rullist):",}
\NormalTok{format: d =\textgreater{} d.label // Shows label in dropdown}
\NormalTok{\})}

\NormalTok{temp = selectedOption.value}
\NormalTok{html\textasciigrave{}\textless{}p\textgreater{}Lagringstemperatur (°C): \textless{}strong\textgreater{}$\{temp.toFixed(0)\}\textless{}/strong\textgreater{}\textless{}/p\textgreater{}}
\NormalTok{\textasciigrave{}}

\NormalTok{viewof pH = Inputs.range([0, 14], \{value: 9,step: 0.1,label: "pH"\})}
\end{Highlighting}
\end{Shaded}

\end{minipage}%
%
\begin{minipage}{0.50\linewidth}

\begin{Shaded}
\begin{Highlighting}[]
\NormalTok{  tillsats= urea*10}
\NormalTok{  tillsatskg = tillsats*va\_vol}
\NormalTok{  sackar = tillsatskg/750}
\NormalTok{  NH3\_NH4 =tillsats*0.46+ammokvave}
\NormalTok{  pKa = (2729.92 / (temp + 273.15)) + 0.090181}
\NormalTok{  f\_nh3 = 1 / (10 ** (pKa {-} pH) + 1)*100}
\NormalTok{  nh3\_gl = f\_nh3/100*NH3\_NH4}
\NormalTok{  nh3\_mm=NH3\_NH4*f\_nh3/14.01*10}
\NormalTok{  dagar = (25{-}temp)+({-}5/({-}0.8*nh3\_gl*0.0037*Math.exp(0.066*temp)))}
  
 

\NormalTok{// Import Observable Plot}
\NormalTok{import \{Plot\} from "@observablehq/plot"}

\NormalTok{// Define your function}
\NormalTok{myFunction = (x) =\textgreater{} (({-}0.8*NH3\_NH4*0.0037*Math.exp(0.066*temp))*x)}

\NormalTok{// Generate data points}
\NormalTok{functionData = \{}
\NormalTok{  const points = [];}
\NormalTok{  for (let x = 0; x \textless{}= 50; x += 0.1) \{}
\NormalTok{    points.push(\{x: x, y: myFunction(x)\});}
\NormalTok{  \}}
\NormalTok{  return points;}
\NormalTok{\}}
\NormalTok{html\textasciigrave{}}
\NormalTok{  \textless{}div class="gray{-}box"\textgreater{}}
\NormalTok{      \textless{}h4\textgreater{}Hygienisering:\textless{}/h4\textgreater{}}
\NormalTok{   Tid för 5{-}log10 reduktion av E. coli:  \textless{}mark\textgreater{}$\{dagar.toFixed(3)\}\textless{}/mark\textgreater{} }
\NormalTok{  dagar\textless{}br\textgreater{}}
\NormalTok{ \textasciigrave{} }


\NormalTok{// Create the plot}
\NormalTok{functionPlot = Plot.plot(\{}
\NormalTok{  title: "Decay",}
\NormalTok{  x: \{}
\NormalTok{    label: "Tid (dagar)", // Add your custom X label here}
\NormalTok{    labelFontSize: 25,     // Make X label bigger}
\NormalTok{    tickSize: 8,}
\NormalTok{    tickFormat: d =\textgreater{} d.toFixed(1), // Format tick numbers}
\NormalTok{    labelOffset: 50  // Reduced offset to keep label closer to axis}
\NormalTok{  \},}
\NormalTok{  y: \{}
\NormalTok{    label: "E. coli (log10/ml)", // Add your custom Y label here}
\NormalTok{    labelFontSize: 12,     // Make Y label bigger}
\NormalTok{    tickSize: 6,}
\NormalTok{    tickFormat: d =\textgreater{} d.toFixed(1), // Format tick numbers}
\NormalTok{    labelOffset: 40  // Add more offset for label to prevent overlap or truncation}
\NormalTok{  \},}
\NormalTok{  width: 600,  // Specify width to have space}
\NormalTok{  height: 450,  // Increased height to accommodate label inside}
\NormalTok{  marginLeft: 80, // Increase left margin to avoid tick label truncation}
\NormalTok{  marginBottom: 80, // Increased bottom margin to keep X label inside plot area}
\NormalTok{  style: \{}
\NormalTok{    fontSize: "22px"       // Make overall plot text bigger}
\NormalTok{  \},}
\NormalTok{  marks: [}
\NormalTok{    Plot.line(functionData, \{x: "x", y: "y", stroke: "\#03865a", strokeWidth: 3\})}
\NormalTok{  ]}
\NormalTok{\})}



\NormalTok{html\textasciigrave{}}
\NormalTok{  \textless{}div class="gray{-}box"\textgreater{}}
\NormalTok{      \textless{}h4\textgreater{}Ureabehandling:\textless{}/h4\textgreater{}}
\NormalTok{     \textless{}p\textgreater{}TillsatsSyrlighet/alkalinitet i rötreste (kg/m3):\textless{}strong\textgreater{} $\{tillsats.toFixed(2)\}\textless{}/strong\textgreater{}\textless{}/p\textgreater{}}
\NormalTok{    \textless{}p\textgreater{}Tillsats (kg): \textless{}strong\textgreater{}$\{tillsatskg.toFixed(0)\}\textless{}/strong\textgreater{}\textless{}/p\textgreater{}}
\NormalTok{    \textless{}p\textgreater{}Antal säckar (à 750 kg): \textless{}strong\textgreater{}$\{sackar.toFixed(1)\}\textless{}/strong\textgreater{}\textless{}/p\textgreater{}}
\NormalTok{    \textless{}p\textgreater{}NH3/NH4+{-}kväve (g/L): \textless{}strong\textgreater{}$\{NH3\_NH4.toFixed(2)\}\textless{}/strong\textgreater{}\textless{}/p\textgreater{}}
\NormalTok{    \textless{}p\textgreater{}pKa: \textless{}strong\textgreater{}$\{pKa.toFixed(2)\}\textless{}/strong\textgreater{}\textless{}/p\textgreater{}}
\NormalTok{    \textless{}p\textgreater{}NH3(\%): \textless{}strong\textgreater{}$\{f\_nh3.toFixed(2)\}\textless{}/strong\textgreater{}\textless{}/p\textgreater{}}
\NormalTok{    \textless{}p\textgreater{}NH3(g/L): \textless{}strong\textgreater{}$\{nh3\_gl.toFixed(2)\}\textless{}/strong\textgreater{}\textless{}/p\textgreater{}}
\NormalTok{    \textless{}p\textgreater{}NH3 (mM): \textless{}strong\textgreater{}$\{nh3\_mm.toFixed(1)\}\textless{}/strong\textgreater{}\textless{}/p\textgreater{}}
\NormalTok{ \textasciigrave{} }
 

\end{Highlighting}
\end{Shaded}

\end{minipage}%
\newline
\begin{minipage}{0.50\linewidth}

\end{minipage}%

\end{figure}%

Tiden
för
nedbrytning
av
urea
beräknas
från
temperaturen
och
är
tillagd
i
den
total
tiden
för
hygienisering.

Modellen
för
inaktivering
av
E.
coli
beroende
på
NH3-koncentration
och
temperatur
(Nordin
2010)
kan
användas
för
både
salmonella
och
E.
coli
inaktivering
då
dessa
har
visat
sig
ha
liknande
inaktivering
vid
ureabehandling.

\begin{center}\rule{0.5\linewidth}{0.5pt}\end{center}

Detta
verktyg
utvecklades
inom
projekt
JTI-21-83-613
finansierat
av
JTI-Stiftelsen
under
tema
Cirkularitet/restströmmar.




\end{document}
